\begin{abstract}
Language models can score proposed survey items cheaply, but their answers mix supplied information with learned associations. We construct synthetic cultural agents from declared aggregate preference anchors. A shared bank of opposing responses receives deterministic labels from the signs of country scores, and Direct Preference Optimization fits a joint choice policy in adapter weights. Evaluation prompts omit country identifiers. In an application linking the Global Preferences Survey to World Values Survey items, the adapters retain trust anchor ordering more strongly than a country-named prompting comparator. Agreement with human country orderings remains uncertain and differs across dimensions. The construction makes the population-specific labeling rule inspectable, while leaving pretrained knowledge and synthetic-bank semantics in the policy. It provides a reusable scoring instrument whose prospective role in survey pre-piloting requires independent human validation.

\medskip
\textcolor{blue}{\textbf{Alternate abstract.} Existing approaches typically generate population-specific responses using country labels or demographic prompts, making it difficult to determine how much of the resulting behavior is driven by the target population profile versus associations encoded during pretraining. We introduce synthetic cultural agents constructed from aggregate preference anchors. Using population-level preference scores, we assign deterministic labels to a bank of contrasting responses and train adapters for different populations with Direct Preference Optimization. Because evaluation prompts contain no population identifiers, the adapters rely on the learned preference profile rather than explicit population cues. We apply this approach to link the General Preference Survey to the World Values Survey. The adapters preserve country-level trust-anchor ordering more accurately than a baseline that names the country in the prompt. Across other preference dimensions, however, the results are mixed. The findings show that aggregate preference-anchored agents may provide a reusable tool for survey-item screening and pre-piloting, but their substantive validity requires independent human evaluation.}

\medskip
\textcolor{purple}{\textbf{Proposed abstract.} Existing approaches typically generate population-specific responses using country labels or demographic prompts, so the resulting behavior mixes a target profile with associations encoded during pretraining. We construct synthetic cultural agents from aggregate preference anchors. Population-level preference scores assign deterministic labels to a shared bank of contrasting responses, and Direct Preference Optimization fits an adapter for each population. Evaluation prompts contain no population identifiers. That omission removes the country-name channel. Pretrained associations and the semantics of the shared bank remain in the fitted policy. We apply the construction to Global Preferences Survey anchors and World Values Survey items. On trust, the adapters retain anchor ordering more strongly than a baseline that names the country in the prompt. Agreement with human country orderings on this panel remains unresolved, and the other preference dimensions are mixed. The labeling rule is inspectable. The fitted policy is a reusable scoring instrument whose prospective role in survey pre-piloting requires independent human validation. Synthetic subjects must not be granted construct validity or permission to speak for people, even when the metrics look favorable.}
\end{abstract}
\noindent\textbf{Keywords:} synthetic cultural agents; aggregate anchors; preference optimization; survey methodology; construct validity

\noindent\textcolor{purple}{\textbf{Keywords (proposed):} synthetic cultural agents; aggregate anchors; preference optimization; survey methodology. [Drop ``construct validity.'' Do not substitute survey pre-piloting. The discussion does not establish that use.]}
